% assets/w05-rotation.svg — 두 오차의 sin·cos 은 회전행렬 하나다
%
%   bash _tools/tikz2svg.sh assets/src/w05-rotation.tex assets/w05-rotation.svg
%
% 요점: 같은 화살표를 북·동 축으로 읽으면 (x-x_k, y-y_k), 경로 축으로 읽으면
%       (x_e, y_e) 다. 둘을 잇는 것이 회전행렬의 전치다. 그리고 t, n 의 꼴은
%       외우는 것이 아니라 방위각의 정의에서 바로 나온다.
%
% 좌표는 계산한 값이다. 화면 규약: TikZ x = 동, TikZ y = 북.
%   pi_p = 50 deg
%     t = (북 cos50, 동 sin50) = (0.642788, 0.766044)   -> TikZ (0.766044, 0.642788)
%     n = (북 -sin50, 동 cos50) = (-0.766044, 0.642788) -> TikZ (0.642788, -0.766044)
%   고른 값: x_e = 42, y_e = 16 (mm)
%     x - x_k = 42 cos50 - 16 sin50 = 26.997 - 12.257 = 14.740   (북)
%     y - y_k = 42 sin50 + 16 cos50 = 32.174 + 10.285 = 42.459   (동)
%     확인: 42.459^2 + 14.740^2 = 2020.0 = 42^2 + 16^2            (길이 보존)
%
%   패널 1 원점 (6,16)   -> 배 (48.46, 30.74)
%   패널 2 원점 (84,24)  -> 배 (126.46, 38.74)
%     t 의 발 = 원점 + 42 t = (116.17, 51.00)   수직성 내적 0.001
%     n 의 발 = 원점 + 16 n = ( 94.29, 11.74)   수직성 내적 0.000
%   패널 3 원점 (18,-48), 단위길이 26 mm
%     t 끝 = (37.92, -31.29),  n 끝 = (34.71, -67.92)
%
% 새로 만든 그림 (2026-10-17). w05-track-error 는 두 오차가 "무엇인지"를 보여 주고,
% 이 그림은 식의 sin·cos 이 "어디서 왔는지"를 보여 준다.

\documentclass[border=6pt]{standalone}
% capstone-style.tex — 이 볼트의 TikZ 그림이 공유하는 스타일
%
% 쓰는 법:  % capstone-style.tex — 이 볼트의 TikZ 그림이 공유하는 스타일
%
% 쓰는 법:  % capstone-style.tex — 이 볼트의 TikZ 그림이 공유하는 스타일
%
% 쓰는 법:  \input{capstone-style}   (assets/src/ 안에서)
% 빌드:     bash _tools/tikz2svg.sh assets/src/w02-node-topic.tex assets/w02-node-topic.svg
%
% 왜 TikZ 인가
%   손으로 SVG 좌표를 쓰면 화살촉·선 굵기·글자 크기가 그림마다 미묘하게 달라진다.
%   그것이 "자료가 어설퍼 보인다"의 정체다. 스타일을 한 번만 정하고 상대좌표로
%   배치하면 좌표를 고쳐도 그림이 무너지지 않는다.
%   수식은 본문과 같은 TeX 글꼴로 나온다.
%
% 한글
%   ko.TeX(DVI 경로)를 쓴다. XeLaTeX 은 TikZ 그래픽을 PDF special 로 내보내는데
%   dvisvgm 이 그것을 받으려면 Ghostscript < 10.01 이 필요하다. 이 PC 에는
%   10.07 이 깔려 있어 그 경로가 막힌다. latex -> DVI -> dvisvgm 이 유일하게
%   동작하는 조합이다.

\usepackage[hangul]{kotex}
\usepackage{tikz}
\usepackage{amsmath}
\usepackage{xcolor}
\usetikzlibrary{arrows.meta, positioning, calc, fit, backgrounds, decorations.pathmorphing}

% ---- 색 --------------------------------------------------------------------
% 강조는 하나만. 나머지는 검정. 흑백 인쇄에서도 읽혀야 한다.
\definecolor{ink}{HTML}{1A1A1A}
\definecolor{soft}{HTML}{666666}
\definecolor{accent}{HTML}{7C3AED}
\definecolor{warn}{HTML}{B91C1C}

% 계층 색 — Simulink 모델의 블록 색과 같은 규칙을 쓴다
\definecolor{guid}{HTML}{1D4ED8}    % 유도
\definecolor{ctrl}{HTML}{EA580C}    % 제어
\definecolor{thr}{HTML}{D97706}     % 추진기
\definecolor{plant}{HTML}{16A34A}   % 운동모델
\definecolor{nav}{HTML}{7C3AED}     % 항법 · 신호처리
\definecolor{pathc}{HTML}{0D9488}   % 계획 경로 · 항적
\definecolor{errc}{HTML}{C2410C}    % 오차
\definecolor{flow}{HTML}{0369A1}    % 조류 · 바람

% ---- 선과 화살촉 -----------------------------------------------------------
\tikzset{
  sig/.style   = {ink, line width=0.5pt, -{Stealth[length=4pt, width=3pt]}},
  wire/.style  = {ink, line width=0.5pt},
  fb/.style    = {ink, line width=0.5pt, densely dashed,
                  -{Stealth[length=4pt, width=3pt]}},
  asig/.style  = {accent, line width=0.5pt, densely dashed,
                  -{Stealth[length=4pt, width=3pt]}},
  % 블록 — 흰 바탕, 직각, 검은 테두리. 색은 테두리에만 준다
  blk/.style   = {draw=ink, line width=0.55pt, fill=white,
                  minimum width=18mm, minimum height=8.5mm, inner sep=2pt, font=\small},
  gblk/.style  = {blk, draw=guid},
  cblk/.style  = {blk, draw=ctrl},
  tblk/.style  = {blk, draw=thr},
  pblk/.style  = {blk, draw=plant},
  nblk/.style  = {blk, draw=nav},
  % 합산점 · 분기점
  sum/.style   = {draw=ink, line width=0.55pt, fill=white, circle,
                  inner sep=0pt, minimum size=4.6mm, font=\scriptsize},
  dot/.style   = {ink, fill=ink, circle, inner sep=0pt, minimum size=1.5mm},
  % 글자
  sname/.style = {font=\small, inner sep=1.5pt},
  note/.style  = {font=\scriptsize, soft, inner sep=1.5pt},
  wnote/.style = {font=\scriptsize, warn, inner sep=1.5pt},
  anote/.style = {font=\scriptsize, accent, inner sep=1.5pt},
  axis/.style  = {ink, line width=0.5pt},
  tick/.style  = {font=\scriptsize, soft},
  % 묶음 상자 — 서브시스템 표시
  grp/.style   = {draw=soft, line width=0.4pt, densely dotted, rounded corners=1mm,
                  inner sep=3mm},
}
   (assets/src/ 안에서)
% 빌드:     bash _tools/tikz2svg.sh assets/src/w02-node-topic.tex assets/w02-node-topic.svg
%
% 왜 TikZ 인가
%   손으로 SVG 좌표를 쓰면 화살촉·선 굵기·글자 크기가 그림마다 미묘하게 달라진다.
%   그것이 "자료가 어설퍼 보인다"의 정체다. 스타일을 한 번만 정하고 상대좌표로
%   배치하면 좌표를 고쳐도 그림이 무너지지 않는다.
%   수식은 본문과 같은 TeX 글꼴로 나온다.
%
% 한글
%   ko.TeX(DVI 경로)를 쓴다. XeLaTeX 은 TikZ 그래픽을 PDF special 로 내보내는데
%   dvisvgm 이 그것을 받으려면 Ghostscript < 10.01 이 필요하다. 이 PC 에는
%   10.07 이 깔려 있어 그 경로가 막힌다. latex -> DVI -> dvisvgm 이 유일하게
%   동작하는 조합이다.

\usepackage[hangul]{kotex}
\usepackage{tikz}
\usepackage{amsmath}
\usepackage{xcolor}
\usetikzlibrary{arrows.meta, positioning, calc, fit, backgrounds, decorations.pathmorphing}

% ---- 색 --------------------------------------------------------------------
% 강조는 하나만. 나머지는 검정. 흑백 인쇄에서도 읽혀야 한다.
\definecolor{ink}{HTML}{1A1A1A}
\definecolor{soft}{HTML}{666666}
\definecolor{accent}{HTML}{7C3AED}
\definecolor{warn}{HTML}{B91C1C}

% 계층 색 — Simulink 모델의 블록 색과 같은 규칙을 쓴다
\definecolor{guid}{HTML}{1D4ED8}    % 유도
\definecolor{ctrl}{HTML}{EA580C}    % 제어
\definecolor{thr}{HTML}{D97706}     % 추진기
\definecolor{plant}{HTML}{16A34A}   % 운동모델
\definecolor{nav}{HTML}{7C3AED}     % 항법 · 신호처리
\definecolor{pathc}{HTML}{0D9488}   % 계획 경로 · 항적
\definecolor{errc}{HTML}{C2410C}    % 오차
\definecolor{flow}{HTML}{0369A1}    % 조류 · 바람

% ---- 선과 화살촉 -----------------------------------------------------------
\tikzset{
  sig/.style   = {ink, line width=0.5pt, -{Stealth[length=4pt, width=3pt]}},
  wire/.style  = {ink, line width=0.5pt},
  fb/.style    = {ink, line width=0.5pt, densely dashed,
                  -{Stealth[length=4pt, width=3pt]}},
  asig/.style  = {accent, line width=0.5pt, densely dashed,
                  -{Stealth[length=4pt, width=3pt]}},
  % 블록 — 흰 바탕, 직각, 검은 테두리. 색은 테두리에만 준다
  blk/.style   = {draw=ink, line width=0.55pt, fill=white,
                  minimum width=18mm, minimum height=8.5mm, inner sep=2pt, font=\small},
  gblk/.style  = {blk, draw=guid},
  cblk/.style  = {blk, draw=ctrl},
  tblk/.style  = {blk, draw=thr},
  pblk/.style  = {blk, draw=plant},
  nblk/.style  = {blk, draw=nav},
  % 합산점 · 분기점
  sum/.style   = {draw=ink, line width=0.55pt, fill=white, circle,
                  inner sep=0pt, minimum size=4.6mm, font=\scriptsize},
  dot/.style   = {ink, fill=ink, circle, inner sep=0pt, minimum size=1.5mm},
  % 글자
  sname/.style = {font=\small, inner sep=1.5pt},
  note/.style  = {font=\scriptsize, soft, inner sep=1.5pt},
  wnote/.style = {font=\scriptsize, warn, inner sep=1.5pt},
  anote/.style = {font=\scriptsize, accent, inner sep=1.5pt},
  axis/.style  = {ink, line width=0.5pt},
  tick/.style  = {font=\scriptsize, soft},
  % 묶음 상자 — 서브시스템 표시
  grp/.style   = {draw=soft, line width=0.4pt, densely dotted, rounded corners=1mm,
                  inner sep=3mm},
}
   (assets/src/ 안에서)
% 빌드:     bash _tools/tikz2svg.sh assets/src/w02-node-topic.tex assets/w02-node-topic.svg
%
% 왜 TikZ 인가
%   손으로 SVG 좌표를 쓰면 화살촉·선 굵기·글자 크기가 그림마다 미묘하게 달라진다.
%   그것이 "자료가 어설퍼 보인다"의 정체다. 스타일을 한 번만 정하고 상대좌표로
%   배치하면 좌표를 고쳐도 그림이 무너지지 않는다.
%   수식은 본문과 같은 TeX 글꼴로 나온다.
%
% 한글
%   ko.TeX(DVI 경로)를 쓴다. XeLaTeX 은 TikZ 그래픽을 PDF special 로 내보내는데
%   dvisvgm 이 그것을 받으려면 Ghostscript < 10.01 이 필요하다. 이 PC 에는
%   10.07 이 깔려 있어 그 경로가 막힌다. latex -> DVI -> dvisvgm 이 유일하게
%   동작하는 조합이다.

\usepackage[hangul]{kotex}
\usepackage{tikz}
\usepackage{amsmath}
\usepackage{xcolor}
\usetikzlibrary{arrows.meta, positioning, calc, fit, backgrounds, decorations.pathmorphing}

% ---- 색 --------------------------------------------------------------------
% 강조는 하나만. 나머지는 검정. 흑백 인쇄에서도 읽혀야 한다.
\definecolor{ink}{HTML}{1A1A1A}
\definecolor{soft}{HTML}{666666}
\definecolor{accent}{HTML}{7C3AED}
\definecolor{warn}{HTML}{B91C1C}

% 계층 색 — Simulink 모델의 블록 색과 같은 규칙을 쓴다
\definecolor{guid}{HTML}{1D4ED8}    % 유도
\definecolor{ctrl}{HTML}{EA580C}    % 제어
\definecolor{thr}{HTML}{D97706}     % 추진기
\definecolor{plant}{HTML}{16A34A}   % 운동모델
\definecolor{nav}{HTML}{7C3AED}     % 항법 · 신호처리
\definecolor{pathc}{HTML}{0D9488}   % 계획 경로 · 항적
\definecolor{errc}{HTML}{C2410C}    % 오차
\definecolor{flow}{HTML}{0369A1}    % 조류 · 바람

% ---- 선과 화살촉 -----------------------------------------------------------
\tikzset{
  sig/.style   = {ink, line width=0.5pt, -{Stealth[length=4pt, width=3pt]}},
  wire/.style  = {ink, line width=0.5pt},
  fb/.style    = {ink, line width=0.5pt, densely dashed,
                  -{Stealth[length=4pt, width=3pt]}},
  asig/.style  = {accent, line width=0.5pt, densely dashed,
                  -{Stealth[length=4pt, width=3pt]}},
  % 블록 — 흰 바탕, 직각, 검은 테두리. 색은 테두리에만 준다
  blk/.style   = {draw=ink, line width=0.55pt, fill=white,
                  minimum width=18mm, minimum height=8.5mm, inner sep=2pt, font=\small},
  gblk/.style  = {blk, draw=guid},
  cblk/.style  = {blk, draw=ctrl},
  tblk/.style  = {blk, draw=thr},
  pblk/.style  = {blk, draw=plant},
  nblk/.style  = {blk, draw=nav},
  % 합산점 · 분기점
  sum/.style   = {draw=ink, line width=0.55pt, fill=white, circle,
                  inner sep=0pt, minimum size=4.6mm, font=\scriptsize},
  dot/.style   = {ink, fill=ink, circle, inner sep=0pt, minimum size=1.5mm},
  % 글자
  sname/.style = {font=\small, inner sep=1.5pt},
  note/.style  = {font=\scriptsize, soft, inner sep=1.5pt},
  wnote/.style = {font=\scriptsize, warn, inner sep=1.5pt},
  anote/.style = {font=\scriptsize, accent, inner sep=1.5pt},
  axis/.style  = {ink, line width=0.5pt},
  tick/.style  = {font=\scriptsize, soft},
  % 묶음 상자 — 서브시스템 표시
  grp/.style   = {draw=soft, line width=0.4pt, densely dotted, rounded corners=1mm,
                  inner sep=3mm},
}


\begin{document}
\begin{tikzpicture}[x=1mm, y=1mm]

% =========================================================================
% 패널 1 — 북·동 축으로 읽기
% =========================================================================
\node[font=\small, anchor=west] at (2,66) {① 북 · 동 축으로 읽으면};

\coordinate (o1) at (6,16);
\coordinate (s1) at (48.46,30.74);

\draw[axis, -{Stealth[length=4pt,width=3pt]}] (o1) -- (6,54);
\draw[axis, -{Stealth[length=4pt,width=3pt]}] (o1) -- (62,16);
\node[font=\scriptsize, anchor=south] at (6,54.6) {$x$ (북)};
\node[font=\scriptsize, anchor=west]  at (62.6,16) {$y$ (동)};

\draw[soft, densely dotted, line width=0.5pt] (s1) -- (48.46,16);
\draw[soft, densely dotted, line width=0.5pt] (s1) -- (6,30.74);
\draw[soft, line width=0.4pt] (46.06,16) -- (46.06,18.4) -- (48.46,18.4);
\draw[soft, line width=0.4pt] (6,28.34) -- (8.4,28.34) -- (8.4,30.74);
\node[font=\scriptsize, soft, anchor=north] at (27,14.6) {$y - y_k$};
\node[font=\scriptsize, soft, anchor=east]  at (4.6,23.4) {$x - x_k$};

\draw[sig, line width=0.8pt] (o1) -- (s1);
\node[font=\scriptsize, anchor=south] at (26.1,26.0) {$\boldsymbol{p} - \boldsymbol{p}_k$};

\fill[pathc] (o1) circle (1.1);
\node[font=\scriptsize, pathc, anchor=north east] at (4.8,14.8) {$\boldsymbol{p}_k$};
\begin{scope}[shift={(48.46,30.74)}, rotate=-20]
  \fill[ink] (-1.6,-2.3) -- (-1.6,1.3) -- (0,3.2) -- (1.6,1.3) -- (1.6,-2.3) -- cycle;
\end{scope}
\node[font=\scriptsize, anchor=west] at (50.2,30.4) {$\boldsymbol{p}$};

\node[note, anchor=west, text width=58mm] at (2,2)
     {축이 북과 동이므로 성분은 좌표를 빼면 바로 나온다. 경로가 어디로 가는지는 들어 있지 않다};

% =========================================================================
% 패널 2 — 경로 축으로 읽기
% =========================================================================
\node[font=\small, anchor=west] at (80,66) {② 경로 축으로 읽으면 — 같은 화살표};

\coordinate (o2)  at (84,24);
\coordinate (s2)  at (126.46,38.74);
\coordinate (ftt) at (116.17,51.00);
\coordinate (ftn) at (94.29,11.74);

% 경로 (t 축) 와 그 오른쪽 (n 축)
\draw[pathc, line width=0.7pt, -{Stealth[length=4pt,width=3pt]}] (o2) -- (126.90,60.00);
\draw[pathc, line width=0.7pt, -{Stealth[length=4pt,width=3pt]}] (o2) -- (96.86,8.68);
\node[font=\scriptsize, pathc, anchor=south east] at (126.6,60.4) {$\hat{t}$ (경로 방향)};
\node[font=\scriptsize, pathc, anchor=north west] at (97.4,8.3)   {$\hat{n}$ (우현)};

% 북 기준선과 경로 방위각
\draw[axis, densely dashed] (o2) -- (84,48);
\draw[axis] (o2) ++(0,14) arc (90:40:14);
\node[note] at (91.2,39.6) {$\pi_p$};

% 성분 — 수선을 내려서 읽는다
\draw[soft, densely dotted, line width=0.5pt] (s2) -- (ftt);
\draw[soft, densely dotted, line width=0.5pt] (s2) -- (ftn);
\draw[soft, line width=0.4pt] (114.33,49.46) -- (115.87,47.62) -- (117.71,49.16);
\draw[soft, line width=0.4pt] (92.75,13.58) -- (94.59,15.12) -- (96.13,13.28);

\draw[errc, line width=1.1pt] (o2) -- (ftt);
\draw[errc, line width=1.1pt] (o2) -- (ftn);
\fill[errc] (ftt) circle (0.9);
\fill[errc] (ftn) circle (0.9);
\node[font=\scriptsize, errc, anchor=south east] at (98.6,39.2) {$x_e$};
\node[font=\scriptsize, errc, anchor=east]       at (87.4,15.4) {$y_e$};

\draw[sig, line width=0.8pt] (o2) -- (s2);
\node[font=\scriptsize, anchor=south] at (104.1,34.0) {$\boldsymbol{p} - \boldsymbol{p}_k$};

\fill[pathc] (o2) circle (1.1);
\node[font=\scriptsize, pathc, anchor=north east] at (82.8,22.8) {$\boldsymbol{p}_k$};
\begin{scope}[shift={(126.46,38.74)}, rotate=-20]
  \fill[ink] (-1.6,-2.3) -- (-1.6,1.3) -- (0,3.2) -- (1.6,1.3) -- (1.6,-2.3) -- cycle;
\end{scope}
\node[font=\scriptsize, anchor=west] at (128.2,38.4) {$\boldsymbol{p}$};

\node[note, anchor=west, text width=58mm] at (80,0)
     {축을 $\pi_p$ 만큼 돌려 놓았을 뿐 화살표는 그대로다 —
      길이도 $\sqrt{x_e^2 + y_e^2}$ 로 변하지 않는다};

% =========================================================================
% 패널 3 — 단위벡터가 왜 그 꼴인가
% =========================================================================
\node[font=\small, anchor=west] at (2,-12) {③ $\hat{t}$ 와 $\hat{n}$ 이 왜 그 꼴인가};

\coordinate (o3) at (18,-48);

\draw[axis, -{Stealth[length=4pt,width=3pt]}] (o3) -- (18,-17);
\draw[axis, -{Stealth[length=4pt,width=3pt]}] (o3) -- (58,-48);
\node[font=\scriptsize, anchor=south west] at (18.6,-17.4) {$x$ (북)};
\node[font=\scriptsize, anchor=west]       at (58.6,-48)   {$y$ (동)};

% 길이 1 을 나타내는 호
\draw[soft, densely dotted, line width=0.4pt] (18,-22) arc (90:-50:26);
\node[note, soft, anchor=west] at (46.19,-37.74) {길이 $1$};

% t 와 그 두 성분
\draw[pathc, line width=1.0pt, -{Stealth[length=5pt,width=4pt]}] (o3) -- (37.92,-31.29);
\node[font=\scriptsize, pathc, anchor=south west] at (37.2,-30.6) {$\hat{t}$};
\draw[pathc, densely dotted, line width=0.5pt] (37.92,-31.29) -- (18,-31.29);
\draw[pathc, densely dotted, line width=0.5pt] (37.92,-31.29) -- (37.92,-48);
\node[font=\scriptsize, pathc, anchor=east]  at (17.4,-33.6) {$\cos\pi_p$};
\node[font=\scriptsize, pathc, anchor=north] at (37.92,-49.2) {$\sin\pi_p$};

% n 과 그 두 성분
\draw[errc, line width=1.0pt, -{Stealth[length=5pt,width=4pt]}] (o3) -- (34.71,-67.92);
\node[font=\scriptsize, errc, anchor=north west] at (34.0,-67.6) {$\hat{n}$};
\draw[errc, densely dotted, line width=0.5pt] (34.71,-67.92) -- (18,-67.92);
\draw[errc, densely dotted, line width=0.5pt] (34.71,-67.92) -- (34.71,-48);
\node[font=\scriptsize, errc, anchor=east]  at (17.4,-67.92) {$-\sin\pi_p$};
\node[font=\scriptsize, errc, anchor=south] at (27,-46.6)    {$\cos\pi_p$};

% 북에서 잰 경로 방위각
\draw[axis] (o3) ++(0,10) arc (90:40:10);
\node[note] at (22.6,-40.6) {$\pi_p$};

% =========================================================================
% 오른쪽 아래 — 식
% =========================================================================
\node[font=\small, anchor=west] at (68,-13) {성분을 읽는 일 = 회전행렬 곱하기};

\node[note, anchor=west, text width=66mm] at (68,-20)
     {$\hat{t}$ 는 방위각이 $\pi_p$ 인 단위벡터 — 북 성분 $\cos\pi_p$, 동 성분 $\sin\pi_p$};
\node[note, anchor=west, text width=66mm] at (68,-29)
     {$\hat{n}$ 은 $\hat{t}$ 를 시계방향으로 $90^\circ$ 돌린 것, 곧 방위각 $\pi_p + 90^\circ$ 인 단위벡터 —
      $\cos(\pi_p{+}90^\circ) = -\sin\pi_p$, $\sin(\pi_p{+}90^\circ) = \cos\pi_p$};
\node[note, anchor=west, text width=66mm] at (68,-40)
     {둘은 길이 1 이고 서로 수직이므로, 성분은 내적 하나로 나온다};

\node[font=\scriptsize, anchor=west] at (68,-48)
     {$x_e = \hat{t}^{\top}(\boldsymbol{p}-\boldsymbol{p}_k)$, \quad
      $y_e = \hat{n}^{\top}(\boldsymbol{p}-\boldsymbol{p}_k)$};

\node[font=\scriptsize, anchor=west] at (68,-59)
     {$\begin{bmatrix} x_e \\[1pt] y_e \end{bmatrix}
       = \begin{bmatrix} \cos\pi_p & \sin\pi_p \\[1pt] -\sin\pi_p & \cos\pi_p \end{bmatrix}
         \begin{bmatrix} x - x_k \\[1pt] y - y_k \end{bmatrix}
       = \boldsymbol{R}(\pi_p)^{\top}\left(\boldsymbol{p}-\boldsymbol{p}_k\right)$};

\node[note, anchor=west, text width=66mm] at (68,-70)
     {$\boldsymbol{R}(\pi_p)$ 는 경로 축의 값을 북 · 동으로 되돌리는 행렬이다.
      반대 방향이므로 전치가 붙는다 — 회전행렬은 $\boldsymbol{R}^{-1} = \boldsymbol{R}^{\top}$ 이라 전치만으로 역이 된다};

\end{tikzpicture}
\end{document}
